\chapter*{Resumen}

La presente tesis desarrolla una guía inteligente contextual para museos orientada a mejorar la experiencia de visita presencial mediante el uso de tecnologías de bajo costo y acceso extendido. El problema abordado radica en que una parte importante de las instituciones museísticas todavía depende de mecanismos de mediación estáticos o poco adaptativos, lo que limita la personalización del recorrido, la accesibilidad y la entrega de información contextual durante la visita. Frente a ello, el objetivo de la investigación es plantear una solución tecnológica viable para reconocer la sala o zona probable en la que se encuentra el visitante, activar contenido pertinente y permitir interacción conversacional apoyada en conocimiento curatorial institucional.

La metodología empleada corresponde a una investigación aplicada con enfoque de diseño de solución. Para ello, se analiza el contexto nacional de adopción tecnológica y demanda museística, se revisa el estado del arte sobre museos inteligentes, localización \textit{indoor}, accesibilidad por voz y recuperación aumentada de información, y se define una arquitectura conceptual integrada. La propuesta combina beacons BLE basados en ESP32, sensores del smartphone para inferencia de orientación, interacción por voz mediante TTS/STT y una capa de inteligencia con RAG para responder consultas con apoyo documental.

Como resultado, la tesis formula una arquitectura funcional y económicamente defendible para mediación cultural digital en el contexto peruano. El aporte principal consiste en articular localización por proximidad, orientación móvil, accesibilidad multimodal e inteligencia conversacional dentro de un enfoque BYOD de bajo costo, escalable y compatible con restricciones operativas de instituciones culturales. Asimismo, la evaluación económico-financiera desarrollada sugiere condiciones favorables de viabilidad para un despliegue institucional progresivo, lo que respalda la pertinencia técnica y práctica de la propuesta.

\noindent\textbf{Palabras clave:} museos inteligentes, BLE, accesibilidad, RAG, mediación cultural digital

